\section{Related work}\label{sec:related}

% DRAFT: every characterisation below is to be checked against the full text of the cited
% work before submission; three works are flagged in the text.

\paragraph{Energy and variational training of neural solvers}
The deep Ritz method \cite{EYu2018} minimises the continuous energy of a variational
problem over neural trial functions, and the deep energy method \cite{Samaniego2020}
brought the idea to computational mechanics. Later work treated essential boundary
conditions \cite{LiaoMing2021} and showed that the numerical quadrature which evaluates the
continuous functional contributes an error of its own, which enters the objective itself
\cite{LiuCaiRamani2023}. Variational and weak-form losses followed: hp-VPINNs
\cite{Kharazmi2021}, conservative energy methods on subdomains \cite{Wang2022CENN},
Kolmogorov--Arnold-informed networks \cite{Wang2025KINN} and energy-based phase-field
fracture \cite{Manav2024}. Robust variational physics-informed networks \cite{Rojas2024}
and the deep Fourier residual method \cite{Taylor2023DFR,Taylor2024} construct losses
equivalent to the energy norm of the error by measuring the residual in a dual norm.
\Cref{lem:exact} and \cref{rem:dual} below are the exact algebraic counterpart of that
equivalence on the assembled system: no quadrature error and no equivalence constants,
because the inverse in the dual norm cancels.

\paragraph{Label-free surrogates and operators}
Variational operator learning \cite{Xu2024VOL} trains operators on the finite-element
system, evaluated matrix-free, and uses the gradient identity of
\cref{lem:exact} as a computational device; its loss is the Euclidean norm of the
residual, or a regression on provisional labels produced by a few conjugate-gradient
iterations. The variational physics-informed neural operator \cite{Eshaghi2025VINO}
evaluates the continuous functional analytically on its own bilinear discretisation over
uniform grids. Other label-free approaches build on the finite-element residual or weak
form \cite{Lee2025FEONet,Uriarte2022,Gao2022,Goswami2022}. Closest to the present work are
graph networks on finite-element meshes trained by minimising the potential energy
\cite{Dalton2023}, physics-informed MeshGraphNets trained without data for non-stationary
and nonlinear problems \cite{Wurth2024}, a physics-informed three-dimensional surrogate for
elastic fields in polycrystals \cite{MonteiroFernandes2025}, and the pretrain finite
element method \cite{Wang2026PretrainFEM}, which pretrains a physics-informed neural
operator without labels and uses its prediction to warm-start the finite-element solver;
its authors report speed-ups of up to an order of magnitude in the warm-start stage. A
follow-up applies that approach to three-dimensional pressure vessels
\cite{Sun2026FEVessel}. Some of these works compare label-free with data-driven training
\cite{Wurth2024,MonteiroFernandes2025}. We are not aware of a comparison of the same
network trained on the same instances with and without their solutions under criteria
fixed in advance, nor of the use of the energy to evaluate and screen surrogate predictions;
both statements, and the characterisations of this paragraph, are to be confirmed against
the full texts of \cite{MonteiroFernandes2025,Wang2026PretrainFEM,Sun2026FEVessel} before
submission.

\paragraph{Supervised mesh surrogates}
MeshGraphNets \cite{Pfaff2021} learn mesh-based simulation by message passing; operator
networks such as DeepONet \cite{Lu2021DeepONet} and the Fourier neural operator
\cite{Li2021FNO} learn maps between function spaces; and point-cloud and
geometry-aware operators handle variable three-dimensional geometries
\cite{He2024GeomDeepONet,Zeng2025PCNO,Liu2026GINOT}. Graph networks have also emulated
cardiac mechanics \cite{Dalton2022}. These surrogates are typically trained and evaluated
with relative $L^2$ errors of the predicted fields. \Cref{rem:rank} and
\cref{sec:results3d} show that the displacement error can rank surrogates in the opposite
order to their stress errors.

\paragraph{Solvers, warm starts and conditioning}
Neural predictions have been used as initial iterates of iterative solvers
\cite{Wang2026PretrainFEM,Eshaghi2026NOWS} and inside hybrid iterative schemes
\cite{Lee2025MetaSolver}. \Cref{cor:warm} bounds how far a warm start can lower the
iteration bound of conjugate gradients, and \cref{sec:results2d:decisions} reports a
pre-registered test in which a warm start saved nothing at a tight tolerance. The training difficulty of
residual-based physics losses has been traced to operator conditioning: descent on a
squared residual is governed by an operator whose condition number is the square of the
original one \cite{DeRyck2024}; the energy objective avoids the squaring
(\cref{rem:dual}).

\paragraph{Fair comparison and pre-registration}
Weak baselines and selective reporting have led to over-optimism in machine learning for
fluid-related partial differential equations \cite{McGreivy2024}, and a comprehensive and
fair comparison of two neural operators needed common data and protocols
\cite{Lu2022Fair}. Pre-registration, the recording of hypotheses, analyses and decision
criteria in a dated document before the data are seen, is established in other sciences as
a guard against these biases \cite{Nosek2018}. Each run whose results decide a claim of
this paper was pre-registered, with executable checks that refuse a run whose configuration
differs from the registered one, and every criterion fixed in advance for these runs is
reported with its outcome (\cref{sec:methods:prereg}, \ref{app:prereg}).
